\documentclass[letterpaper,twocolumn,10pt]{article}
\usepackage{usenix_style}
\usepackage{amsmath,amssymb,amsfonts}
\usepackage{graphicx}
\usepackage{textcomp}
\usepackage{amsthm}
\newtheorem{definition}{Definition}

\begin{document}

\date{}

\title{\Large \bf Nắm Bắt Chiến Dịch, Không Phải Kết Nối: Mạng Siêu Đồ Thị Nhận Thức Phối Hợp Cho Phát Hiện Xâm Nhập Đa Giai Đoạn}

\author{
{\rm Anonymous Authors}\\
Institution
}
\maketitle

\begin{abstract}
Hệ thống phát hiện xâm nhập dựa trên nguồn gốc (provenance) phân tích nhật ký kiểm toán (audit logs) cấp hạt nhân để xác định các Mối đe dọa dai dẳng nâng cao (APT), nhưng các phương pháp hiện tại hoạt động ở mức độ chi tiết thô --- phân loại toàn bộ đồ thị (KAIROS, MAGIC) hoặc các nút riêng lẻ (FLASH, ThreaTrace) --- mà không quy kết việc phát hiện cho các sự kiện tấn công cụ thể. Chúng tôi giới thiệu HyperMamba, một hệ thống phát hiện các sự kiện tấn công xuyên tiến trình (crossprocess) riêng lẻ trong các luồng kiểm toán thời gian thực bằng cách duy trì trạng thái thực thể liên tục thông qua Mô hình Không gian Trạng thái Chọn lọc (Selective State Space Models). HyperMamba mô hình hóa mỗi sự kiện kiểm toán như một siêu cạnh (hyperedge) thời gian kết nối hai hoặc ba thực thể hệ thống, bảo tồn cấu trúc biến thiên nguyên bản của Mô hình Dữ liệu Chung (CDM). Một SSM Chọn lọc duy trì một vectơ trạng thái cho mỗi thực thể, tích lũy bối cảnh hành vi trên toàn bộ luồng; tại mỗi sự kiện, tập hợp siêu cạnh dựa trên cơ chế chú ý (attention) kết hợp trạng thái của các thực thể tham gia, SSM cập nhật mỗi thực thể và một bộ phân loại tạo ra quyết định phát hiện ở cấp độ sự kiện. Chúng tôi chứng minh bằng thực nghiệm rằng trạng thái tích lũy này là cơ chế phát hiện chính: một phép cắt bỏ (ablation) loại bỏ trạng thái liên tục trong khi vẫn giữ nguyên bộ mã hóa sự kiện và bộ phân loại làm giảm AUPRC kiểm tra từ 0.73 xuống 0.33. Một cơ chế điều biến danh tính tiến trình có cổng (gated process identity modulation) cung cấp các điểm neo ngữ nghĩa bất biến theo chiến dịch, cho phép tổng quát hóa sang các chiến dịch tấn công hoàn toàn chưa từng thấy. Trên tập dữ liệu DARPA TC Engagement 3 Theia, HyperMamba đạt được AUPRC cấp sự kiện là 0.86 trong điều kiện tiếp tục cùng chiến dịch (same-campaign) và 0.72 trong điều kiện chéo chiến dịch (cross-campaign), với tỷ lệ dương tính giả lần lượt là 0.000010 và 0.000195, trong khi xử lý các sự kiện với tốc độ gấp $194\times$ tốc độ thu thập dữ liệu cao nhất trên một GPU duy nhất.
\end{abstract}

\section{Giới thiệu}
\label{sec:introduction}

Hệ thống Phát hiện Xâm nhập Mạng (NIDS) dựa trên nguồn gốc phân tích nhật ký kiểm toán cấp hạt nhân để phát hiện các Mối đe dọa dai dẳng nâng cao (APT) trên các máy chủ được giám sát. Các hệ thống tiên tiến gần đây --- KAIROS~\cite{kairos}, MAGIC~\cite{magic}, FLASH~\cite{flash} và ThreaTrace~\cite{threatrace} --- đã thể hiện khả năng phát hiện mạnh mẽ. Tuy nhiên, các hệ thống này hoạt động ở mức độ chi tiết về cơ bản là thô: KAIROS và MAGIC gán điểm bất thường cho toàn bộ đồ thị nguồn gốc, trong khi FLASH và ThreaTrace phân loại các nút (tiến trình hoặc tệp) riêng lẻ. Không có hệ thống nào đưa ra phân loại ở \emph{cấp độ sự kiện}. Khi một nhà phân tích bảo mật nhận được cảnh báo cấp nút đánh dấu một tiến trình là độc hại, nhà phân tích vẫn phải kiểm tra thủ công hàng ngàn sự kiện của tiến trình đó để xác định hoạt động cụ thể nào cấu thành cuộc tấn công. Khoảng trống pháp y giữa phát hiện và quy kết này làm hạn chế tiện ích hoạt động của các hệ thống hiện tại.

Khoảng trống chi tiết không chỉ đơn thuần là bất tiện --- nó phản ánh một hạn chế về cấu trúc. Các sự kiện tấn công xuyên tiến trình, được thực thi bởi các tiến trình con sinh ra trong các chiến dịch APT đa giai đoạn, sử dụng các lệnh gọi hệ thống (system calls) giống hệt như hoạt động nền lành tính. Phân phối đặc trưng trên mỗi sự kiện của chúng thỏa mãn $P(\mathbf{x}_i \mid \text{tấn công}) \approx P(\mathbf{x}_i \mid \text{lành tính})$, khiến chúng không thể phân biệt được về mặt thống kê nếu xét riêng lẻ. Một hệ thống đánh giá các sự kiện một cách độc lập không thể tách biệt một cách đáng tin cậy các hoạt động đọc và ghi độc hại khỏi các hoạt động lành tính. Tín hiệu phân biệt chỉ tồn tại trong \emph{lịch sử hành vi tích lũy} của thực thể thực thi: chuỗi các hoạt động trước đó tiết lộ vai trò của thực thể trong một chiến dịch đang diễn ra. Các hệ thống hiện tại loại bỏ tín hiệu này bằng cách hoạt động trên các ảnh chụp nhanh (snapshots) đồ thị tĩnh mà không có trạng thái thực thể liên tục.

Chúng tôi giới thiệu \textbf{HyperMamba}, một NIDS dựa trên nguồn gốc giúp phát hiện các sự kiện tấn công xuyên tiến trình riêng lẻ trong các luồng kiểm toán thời gian thực bằng cách duy trì trạng thái thực thể liên tục thông qua Mô hình Không gian Trạng thái Chọn lọc. HyperMamba mô hình hóa mỗi sự kiện kiểm toán như một siêu cạnh thời gian kết nối hai hoặc ba thực thể hệ thống, bảo tồn cấu trúc biến thiên nguyên bản của Mô hình Dữ liệu Chung (CDM). Một SSM Chọn lọc duy trì và cập nhật một vectơ trạng thái cho mỗi thực thể trong hệ thống, tích lũy bối cảnh hành vi trên toàn bộ luồng sự kiện. Tại mỗi sự kiện, mô hình tập hợp trạng thái của các thực thể tham gia thông qua tập hợp siêu cạnh dựa trên chú ý, cập nhật trạng thái của từng thực thể thông qua vòng lặp SSM và phân loại sự kiện bằng cách sử dụng cả trạng thái được cập nhật và các đặc trưng của chính sự kiện đó. Kết quả là khả năng phát hiện thời gian thực, cấp độ sự kiện với quy kết pháp y cho mỗi sự kiện.

Những đóng góp của chúng tôi bao gồm:
\begin{enumerate}
    \item \textbf{Hình thức hóa siêu đồ thị nguồn gốc (C1).} Chúng tôi hình thức hóa các luồng kiểm toán thành các siêu đồ thị thời gian với siêu cạnh có kích thước biến thiên và chứng minh thông qua phép cắt bỏ có kiểm soát rằng việc tập hợp trạng thái xuyên thực thể trong các siêu cạnh mang lại sự cải thiện - AUPRC so với việc theo dõi từng thực thể riêng biệt.

    \item \textbf{Trạng thái thực thể liên tục làm cơ chế phát hiện (C2).} Chúng tôi chứng minh bằng thực nghiệm rằng trạng thái thực thể tích lũy --- không phải các đặc trưng sự kiện riêng biệt --- là cơ chế chính cho phép phát hiện tấn công xuyên tiến trình. Loại bỏ trạng thái liên tục làm sụp đổ AUPRC kiểm tra từ 0.73 xuống 0.33 (Mục~\ref{sec:ablation_state}).

    \item \textbf{Danh tính tiến trình có cổng cho tổng quát hóa chéo chiến dịch (C3).} Chúng tôi giới thiệu một cơ chế điều biến danh tính tiến trình có cổng cung cấp các điểm neo ngữ nghĩa bất biến theo chiến dịch, cho phép mô hình phát hiện tổng quát hóa cho các chiến dịch tấn công hoàn toàn chưa từng thấy (Mục~\ref{sec:ablation_process}).

    \item \textbf{Phát hiện cấp sự kiện ở quy mô hoạt động (C4).} HyperMamba đạt AUPRC cấp sự kiện là 0.86 (cùng chiến dịch) và 0.72 (chéo chiến dịch) với FPR lần lượt là $0.000010$ và $0.000195$, trong khi xử lý gấp $194\times$ tốc độ thu thập dữ liệu kiểm toán đỉnh điểm trên một GPU duy nhất (Mục~\ref{sec:main_results}).
\end{enumerate}

Mục~\ref{sec:background} giới thiệu mô hình dữ liệu nguồn gốc và kiến thức nền tảng về SSM. Mục~\ref{sec:threat_model} định nghĩa mô hình mối đe dọa và bài toán phát hiện. Mục~\ref{sec:architecture} trình bày chi tiết kiến trúc HyperMamba. Mục~\ref{sec:evaluation} trình bày đánh giá. Mục~\ref{sec:related_work} và~\ref{sec:discussion} thảo luận về các công trình liên quan và những hạn chế.

\section{Kiến thức Nền tảng}
\label{sec:background}

Mục này định nghĩa ba khái niệm cần thiết để theo dõi kiến trúc (Mục~\ref{sec:architecture}): mô hình dữ liệu nguồn gốc, sự hình thức hóa của nó dưới dạng siêu đồ thị thời gian, và vòng lặp của Mô hình Không gian Trạng thái Chọn lọc.

\subsection{Mô hình Dữ liệu Nguồn gốc}
\label{sec:provenance}

Các hệ thống kiểm toán cấp máy chủ hiện đại dựa trên Mô hình Dữ liệu Chung (CDM) của DARPA ghi lại mọi hoạt động qua trung gian hạt nhân dưới dạng một sự kiện có kiểu kết nối một \emph{chủ thể} (tiến trình thực thi) với một hoặc hai \emph{đối tượng} (mục tiêu của hoạt động). Một thao tác đọc tệp tạo ra một sự kiện với một chủ thể và một đối tượng (\texttt{tiến trình} $\to$ \texttt{tệp}); một thao tác ánh xạ bộ nhớ của một thư viện chia sẻ tạo ra một sự kiện với một chủ thể và hai đối tượng (\texttt{tiến trình} $\to$ \texttt{tệp} $\to$ \texttt{đối tượng\_bộ\_nhớ}). Mỗi sự kiện mang siêu dữ liệu: một loại sự kiện rời rạc (một trong $\sim$40 loại CDM như \texttt{EVENT\_READ}, \texttt{EVENT\_MMAP}), các thuộc tính liên tục (số byte, cờ), và dấu thời gian với độ phân giải nano giây. Trong một chiến dịch kéo dài nhiều ngày, một máy chủ duy nhất có thể tạo ra hàng chục triệu sự kiện như vậy.

\subsection{Siêu đồ thị Nguồn gốc}
\label{sec:hypergraph}

Chúng tôi hình thức hóa luồng kiểm toán dưới dạng một siêu đồ thị thời gian.

\begin{definition}[Siêu đồ thị Nguồn gốc]
\label{def:hypergraph}
Một \emph{siêu đồ thị nguồn gốc} là một bộ $\mathcal{G} = (\mathcal{V}, \mathcal{E})$, trong đó $\mathcal{V}$ là một tập hữu hạn các thực thể hệ thống (tiến trình, tệp, socket, đối tượng bộ nhớ) và $\mathcal{E} = \{e_1, e_2, \dots\}$ là một chuỗi các siêu cạnh được sắp xếp theo thời gian. Mỗi siêu cạnh $e_t = (v_{\text{subj}}, v_{\text{obj}}, v_{\text{obj2}}, \tau_t, \mathbf{a}_t)$ kết nối $|V_{e_t}| \in \{2, 3\}$ thực thể, trong đó $v_{\text{obj2}}$ vắng mặt đối với các sự kiện kích thước 2, $\tau_t$ là dấu thời gian của sự kiện, và $\mathbf{a}_t$ là vectơ thuộc tính.
\end{definition}

\noindent Ví dụ, khi tiến trình Firefox bị xâm phạm ánh xạ tệp nhị phân độc hại \texttt{malicious.so} vào bộ nhớ có thể thực thi, siêu cạnh \texttt{EVENT\_MMAP} được tạo ra kết nối ba thực thể: Firefox ($v_{\text{subj}}$), tệp \texttt{malicious.so} ($v_{\text{obj}}$), và vùng bộ nhớ được cấp phát ($v_{\text{obj2}}$).

\textbf{Tại sao lại là siêu cạnh, không phải cạnh cặp đôi.} Các NIDS dựa trên nguồn gốc trước đây~\cite{kairos,magic,flash,threatrace} phân rã mỗi sự kiện đa đối tượng thành các cạnh cặp đôi độc lập, loại bỏ cấu trúc đồng xuất hiện. Trong tập dữ liệu DARPA TC E3 Theia, 2.8\% các sự kiện là siêu cạnh kích thước 3. Mặc dù tỷ lệ này có vẻ nhỏ, các sự kiện kích thước 3 lại xuất hiện nhiều gấp $3.0\times$ trong số các sự kiện tấn công xuyên tiến trình so với các sự kiện lành tính, chủ yếu do các hoạt động \texttt{EVENT\_MMAP} (ví dụ: một tiến trình bị xâm phạm ánh xạ một thư viện chia sẻ độc hại vào bộ nhớ thực thi). Sự làm giàu này là lý do để duy trì kích thước (arity) nguyên bản của siêu cạnh thay vì phân rã thành các tương tác cặp đôi.

\subsection{Mô hình Không gian Trạng thái Chọn lọc}
\label{sec:ssm_background}

Mô hình Không gian Trạng thái (SSMs) xác định một phép đệ quy tuyến tính trên một trạng thái ẩn $\mathbf{h} \in \mathbb{R}^n$, ánh xạ một chuỗi đầu vào $x(t)$ tới đầu ra $y(t)$ thông qua động lực học thời gian liên tục:
\begin{equation}
    \mathbf{h}'(t) = \mathbf{A}\mathbf{h}(t) + \mathbf{B}x(t), \quad y(t) = \mathbf{C}\mathbf{h}(t)
\end{equation}
trong đó $\mathbf{A} \in \mathbb{R}^{n \times n}$ điều chỉnh sự suy giảm trạng thái và $\mathbf{B} \in \mathbb{R}^{n \times 1}$ kiểm soát việc tiêm đầu vào. Đối với các chuỗi rời rạc, động lực học liên tục được rời rạc hóa với bước nhảy $\Delta$ để mang lại:
\begin{equation}
    \bar{\mathbf{A}} = \exp(\Delta \mathbf{A}), \quad \bar{\mathbf{B}} = (\Delta \mathbf{A})^{-1}(\bar{\mathbf{A}} - \mathbf{I}) \cdot \Delta \mathbf{B}
\end{equation}
\begin{equation}
    \mathbf{h}_t = \bar{\mathbf{A}} \mathbf{h}_{t-1} + \bar{\mathbf{B}} x_t, \quad y_t = \mathbf{C} \mathbf{h}_t
\end{equation}

SSM \emph{Chọn lọc}~\cite{mamba} khiến $\mathbf{B}$, $\mathbf{C}$, và $\Delta$ phụ thuộc vào đầu vào, cho phép mô hình chọn lọc giữ lại hoặc loại bỏ thông tin tại mỗi bước dựa trên đầu vào hiện tại. Trong Mục~\ref{sec:ssm_update}, chúng tôi điều chỉnh cơ chế này để hoạt động trên các trạng thái thực thể bên trong một siêu đồ thị, nơi $\Delta$ được điều kiện hóa thêm dựa trên thời gian trôi qua kể từ lần cuối cùng thực thể được quan sát.

\section{Mô hình Mối đe dọa}
\label{sec:threat_model}

\subsection{Mô hình Kẻ tấn công}
\label{sec:attacker_model}

Chúng tôi xem xét một kẻ tấn công thuộc nhóm Mối đe dọa dai dẳng nâng cao (APT) đang thực thi các chiến dịch tấn công xuyên tiến trình đa giai đoạn trên một máy chủ được giám sát.

\textbf{Kiến thức.} Kẻ tấn công biết rằng tính năng ghi nhật ký kiểm toán cấp hạt nhân đang hoạt động trên máy chủ đích. Trong đánh giá chính, kẻ tấn công không có kiến thức về kiến trúc hoặc các tham số của mô hình phát hiện (cài đặt hộp đen). Chúng tôi thảo luận về các tác động của một kẻ tấn công hộp trắng nhắm mục tiêu vào cơ chế truyền bá trạng thái SSM trong Mục~\ref{sec:discussion}.

\textbf{Khả năng.} Kẻ tấn công có thể thực thi các hoạt động tùy ý thông qua các API hệ thống tiêu chuẩn: sinh ra các tiến trình, đọc và ghi tệp, ánh xạ các vùng bộ nhớ và thiết lập kết nối mạng. Tất cả các hoạt động như vậy đều được cơ sở hạ tầng ghi nhật ký của hạt nhân ghi lại dưới dạng các sự kiện kiểm toán. Kẻ tấn công \emph{không thể} can thiệp hoặc chặn nhật ký kiểm toán cấp hạt nhân. Giả định này là tiêu chuẩn trong các tài liệu NIDS dựa trên nguồn gốc và phản ánh các hệ thống triển khai nơi việc thu thập kiểm toán chạy ở mức đặc quyền cao hơn các tiến trình không gian người dùng.

\textbf{Mục tiêu.} Mục tiêu của kẻ tấn công là thực thi một chiến dịch xuyên tiến trình đa giai đoạn hoàn chỉnh --- khai thác ban đầu một ứng dụng điểm vào, triển khai các công cụ thứ cấp thông qua các tiến trình con, thiết lập các kênh điều khiển và chỉ huy (C2), và rò rỉ dữ liệu --- trong khi trốn tránh sự phát hiện ở cấp độ sự kiện. Kẻ tấn công đạt được sự trốn tránh nếu các sự kiện hệ thống riêng lẻ do các tiến trình con của chiến dịch tạo ra không bị hệ thống phát hiện đánh dấu là độc hại.

\subsection{Sự kiện Tấn công DARPA TC E3}
\label{sec:attack_instantiation}

Chúng tôi dựa trên mô hình mối đe dọa vào một chiến dịch đa giai đoạn cụ thể từ tập dữ liệu DARPA TC Engagement 3 Theia. Cuộc tấn công diễn ra qua bốn giai đoạn riêng biệt, mỗi giai đoạn tạo ra các sự kiện kiểm toán trong nhật ký cấp hạt nhân của máy chủ:

\begin{enumerate}
    \item \textbf{Xâm phạm ban đầu.} Kẻ tấn công khai thác một lỗ hổng trong tiến trình trình duyệt Firefox. Firefox thực thi một \texttt{EVENT\_WRITE} để thả một mã nhị phân độc hại tại \texttt{/home/admin/clean}. Sự kiện đơn lẻ này là \emph{điểm vào} --- Firefox là tiến trình bị xâm phạm.

    \item \textbf{Kích hoạt tải trọng.} Firefox sinh ra tiến trình \texttt{clean} thông qua \texttt{EVENT\_EXECUTE}. Tệp nhị phân \texttt{clean} tải các thư viện chia sẻ (\texttt{EVENT\_MMAP}) và đọc các tệp cấu hình hệ thống (\texttt{EVENT\_READ} trên \texttt{/etc/resolv.conf}, \texttt{/etc/hosts}).

    \item \textbf{Điều khiển và Chỉ huy (C2).} \texttt{clean} thiết lập một kết nối reverse shell đến cơ sở hạ tầng của kẻ tấn công thông qua \texttt{EVENT\_CONNECT} đến một địa chỉ IP bên ngoài. Các sự kiện \texttt{EVENT\_SENDTO} và \texttt{EVENT\_RECVFROM} tiếp theo mang các lệnh và phản hồi C2.

    \item \textbf{Rò rỉ dữ liệu.} Dưới sự điều hướng của C2, \texttt{clean} đọc các tệp nhạy cảm (\texttt{EVENT\_READ}) và truyền nội dung của chúng tới máy chủ của kẻ tấn công (\texttt{EVENT\_SENDTO}). Các tiến trình con bổ sung có thể được sinh ra bởi \texttt{clean} để thực hiện việc trinh sát mạng ngang hàng.
\end{enumerate}

\textbf{Định nghĩa nhãn xuyên tiến trình (crossprocess).} Trong bài báo này, nhãn \emph{crossprocess} đề cập đến các sự kiện được thực thi bởi các tiến trình con do tiến trình điểm vào sinh ra --- trong trường hợp này, \texttt{clean} và bất kỳ tiến trình nào mà nó sinh ra sau đó. Bản thân tiến trình điểm vào (Firefox) bị \emph{loại trừ} khỏi nhãn này. Việc loại trừ này là có chủ ý: việc khai thác ban đầu của Firefox có thể tạo ra các lệnh gọi hệ thống bất thường (ví dụ: trình duyệt ghi một tệp thực thi), có thể phát hiện được bằng các phương pháp không trạng thái, đơn giản hơn. Bài toán phát hiện khó hơn --- và là mục tiêu của HyperMamba --- là xác định các tiến trình con ở giai đoạn sau, những sự kiện riêng lẻ của chúng về mặt vận hành không thể phân biệt được với hoạt động hệ thống lành tính.

\subsection{Tại sao Tấn công Xuyên tiến trình lại khó đối với Phát hiện Không trạng thái}
\label{sec:crossprocess_hard}

Các cuộc tấn công xuyên tiến trình được thực thi bởi \emph{các tiến trình con} mà tiến trình điểm vào bị xâm phạm sinh ra để thực hiện các giai đoạn hoạt động của chiến dịch --- trinh sát, di chuyển ngang, tập hợp dữ liệu và rò rỉ. Các tiến trình con này (ví dụ: mã nhị phân \texttt{clean} hoặc một tập lệnh shell được tiêm vào) sử dụng chính xác các lệnh gọi hệ thống giống như các tiến trình nền lành tính: \texttt{EVENT\_READ} trên tệp cấu hình, \texttt{EVENT\_WRITE} vào các thư mục tạm thời, \texttt{EVENT\_MMAP} để tải thư viện chia sẻ, và \texttt{EVENT\_CONNECT} đến các máy chủ bên ngoài. Một bộ phân loại không trạng thái, đánh giá mỗi sự kiện một cách riêng lẻ, quan sát các lệnh gọi này mà không có bối cảnh lịch sử và đối mặt với một sự mơ hồ cơ bản.

\textbf{Ví dụ cụ thể.} Hãy xem xét hai sự kiện được rút ra từ tập dữ liệu DARPA TC E3 Theia:
\begin{enumerate}
    \item \texttt{bash} (lành tính) thực thi \texttt{EVENT\_READ} trên \texttt{/etc/resolv.conf}, chuyển 242 byte, $\Delta t = 0.4$s kể từ sự kiện cuối cùng của nó.
    \item \texttt{clean} (tấn công xuyên tiến trình) thực thi \texttt{EVENT\_READ} trên \texttt{/etc/resolv.conf}, chuyển 242 byte, $\Delta t = 0.3$s kể từ sự kiện cuối cùng của nó.
\end{enumerate}
Cả hai sự kiện chia sẻ cùng loại sự kiện, cùng tệp đích và các đặc trưng liên tục (số byte, cờ, thời gian giữa các lần đến) gần như giống hệt nhau. Một vectơ đặc trưng được trích xuất từ bất kỳ sự kiện nào một cách riêng lẻ đều không thể phân biệt được về mặt thống kê. Một cách chính thức hơn, các phân phối đặc trưng trên mỗi sự kiện thỏa mãn $P(\mathbf{x}_i \mid \text{tấn công}) \approx P(\mathbf{x}_i \mid \text{lành tính})$ đối với phần lớn các sự kiện tấn công xuyên tiến trình. Tín hiệu phân biệt không nằm ở bản thân sự kiện, mà nằm ở \emph{lịch sử hành vi tích lũy} của thực thể chủ thể: \texttt{clean} được sinh ra bởi tiến trình Firefox bị xâm phạm, trước đó đã thiết lập kết nối C2 và đã ghi vào các thư mục tập hợp dữ liệu --- một mô hình chỉ có thể nhìn thấy thông qua trạng thái liên tục.

\textbf{Kiểm chứng thực nghiệm.} Nghiên cứu cắt bỏ trong Mục~\ref{sec:ablation_state} cung cấp bằng chứng thực nghiệm trực tiếp cho điều kiện lẩn tránh này. Việc loại bỏ trạng thái thực thể tích lũy khỏi HyperMamba (trong khi vẫn giữ lại toàn bộ bộ mã hóa sự kiện, bộ tập hợp siêu cạnh và bộ phân loại) khiến AUPRC kiểm tra sụp đổ từ 0.73 xuống 0.33. AUPRC đào tạo cũng giảm từ 0.86 xuống 0.51, xác nhận rằng mô hình thậm chí không thể khớp với phân phối đào tạo nếu chỉ sử dụng các đặc trưng sự kiện cục bộ. Các sự kiện tấn công xuyên tiến trình, theo cấu trúc, không thể xác định được nếu không có bối cảnh tích lũy.

\subsection{Bài toán Phát hiện}
\label{sec:detection_task}

Cho một luồng sự kiện kiểm toán liên tục theo trình tự thời gian $e_1, e_2, \dots$ được tạo ra bởi một máy chủ giám sát duy nhất, bài toán phát hiện là phân loại mỗi sự kiện $e_t$ là lành tính ($y_t = 0$) hoặc tấn công xuyên tiến trình ($y_t = 1$) tại thời điểm sự kiện đến. Về mặt hình thức, bộ phát hiện là một hàm:
\begin{equation}
    f_\theta(e_t \mid e_1, \dots, e_{t-1}) \to [0, 1]
\end{equation}
tạo ra xác suất $\hat{y}_t$ sử dụng sự kiện hiện tại $e_t$ làm đầu vào và tất cả các sự kiện đã quan sát trước đó $e_1, \dots, e_{t-1}$ làm bối cảnh lịch sử. Bộ phát hiện \emph{không thể} truy cập các sự kiện tương lai $e_{t+1}, e_{t+2}, \dots$ --- việc phân loại phải mang tính nhân quả nghiêm ngặt.

\textbf{Phạm vi.} Nghiên cứu này giải quyết việc phát hiện dựa trên máy chủ từ nhật ký kiểm toán cấp hạt nhân. Phát hiện xâm nhập dựa trên luồng mạng, kiểm tra tải trọng, và khớp chữ ký nằm ngoài phạm vi của bài báo này.

\section{Kiến trúc HyperMamba}
\label{sec:architecture}

HyperMamba xử lý các sự kiện kiểm toán theo trình tự thời gian nghiêm ngặt dưới dạng luồng liên tục. Mỗi sự kiện là một siêu cạnh thời gian kết nối hai hoặc ba thực thể hệ thống (Mục~\ref{sec:hypergraph}). Với mỗi sự kiện đến, mô hình thực hiện một lượt truyền thẳng (forward pass) năm giai đoạn để đọc và cập nhật các trạng thái thực thể liên tục. Hình~\ref{fig:architecture} tóm tắt luồng xử lý; các Mục~\ref{sec:event_encoding}--\ref{sec:training} trình bày chi tiết từng giai đoạn.

\subsection{Tổng quan}
\label{sec:overview}

Gọi $e_t$ là sự kiện kiểm toán đến tại bước rời rạc $t$, với các thực thể tham gia $V_{e_t} \subseteq \mathcal{V}$, trong đó $|V_{e_t}| \in \{2, 3\}$. Chúng tôi ký hiệu ba khe thực thể là $v_{\text{subj}}$ (tiến trình chủ thể), $v_{\text{obj}}$ (đối tượng chính), và $v_{\text{obj2}}$ (đối tượng phụ, vắng mặt với sự kiện kích thước 2). Mỗi thực thể $v \in \mathcal{V}$ duy trì một vectơ trạng thái liên tục $\mathbf{s}_v \in \mathbb{R}^d$ trong \texttt{EntityStateBank}, với $d = 128$ trong suốt bài báo này.

HyperMamba xử lý $e_t$ theo năm giai đoạn:

\begin{enumerate}
    \item \textbf{Mã hóa Sự kiện} (Mục~\ref{sec:event_encoding}). Loại sự kiện, các đặc trưng liên tục (số byte, cờ, thống kê thời gian đến giữa các sự kiện), và đường dẫn tiến trình của chủ thể được chiếu và kết hợp thành biểu diễn sự kiện dày đặc $\mathbf{f}_{e} \in \mathbb{R}^{2d}$. Danh tính đường dẫn tiến trình được tiêm thông qua cổng thặng dư cộng với dropout danh tính ở giai đoạn đào tạo để ngăn việc ghi nhớ theo từng tiến trình (Mục~\ref{sec:event_encoding}).
    
    \item \textbf{Truy xuất Trạng thái Thực thể} (Mục~\ref{sec:entity_bank}). Đối với mỗi thực thể $v \in V_{e_t}$, trạng thái hiện tại $\mathbf{s}_v$ và thời gian đã trôi qua $\Delta t_v$ kể từ lần cuối cùng $v$ được quan sát được đọc từ \texttt{EntityStateBank}. Một mặt nạ hợp lệ $m_v$ đánh dấu các đối tượng phụ vắng mặt.
    
    \item \textbf{Tập hợp Siêu cạnh}, $\mathcal{V} \to \mathcal{E}$ (Mục~\ref{sec:aggregation}). Các trạng thái thực thể được truy xuất sẽ được tập hợp thành một biểu diễn siêu cạnh duy nhất $\mathbf{x}_e \in \mathbb{R}^d$ thông qua cơ chế chú ý tích vô hướng có chia tỷ lệ, nơi các đặc trưng sự kiện $\mathbf{f}_e$ đóng vai trò là truy vấn và các trạng thái thực thể --- được nối với các nhúng vai trò học được $\mathbf{r}_v$ và thời gian trôi qua theo tỷ lệ logarit $\log(1 + \Delta t_v)$ --- đóng vai trò là khóa và giá trị.
    
    \item \textbf{Cập nhật Trạng thái SSM}, $\mathcal{E} \to \mathcal{V}$ (Mục~\ref{sec:ssm_update}). Trạng thái của mỗi thực thể tham gia $\mathbf{s}_v$ được cập nhật thông qua Mô hình Không gian Trạng thái Chọn lọc. Cập nhật này đặc thù theo vai trò (chủ thể và đối tượng nhận cập nhật bất đối xứng), nhận biết thời gian (bước rời rạc $\Delta_i$ tích hợp $\Delta t_v$), và có cổng (một cổng học được kiểm soát lượng đầu ra SSM sẽ thay thế trạng thái trước đó).
    
    \item \textbf{Phân loại} (Mục~\ref{sec:classification}). Các trạng thái thực thể sau khi cập nhật được lấy trung bình (được che bằng mặt nạ hợp lệ) và nối với $\mathbf{f}_e$ để tạo ra một logit vô hướng $\hat{y}_t = \sigma(\text{MLP}([\bar{\mathbf{s}} \parallel \mathbf{W}_p \mathbf{f}_e]))$, trong đó $\hat{y}_t \in [0, 1]$ là xác suất $e_t$ là một sự kiện tấn công xuyên tiến trình.
\end{enumerate}

Trạng thái được ghi vào ngân hàng ở giai đoạn 4 cho thực thể $v$ chính là trạng thái được đọc ở giai đoạn 2 lần tiếp theo $v$ tham gia vào bất kỳ sự kiện nào. Điều này tạo ra một đường dẫn thông tin hồi quy giữa các thực thể và qua thời gian, với bộ nhớ giới hạn ở $\mathcal{O}(|\mathcal{V}| \times d)$ không phụ thuộc vào số lượng sự kiện đã xử lý. Việc đào tạo sử dụng Lan truyền ngược qua thời gian bị cắt ngắn (Truncated Backpropagation Through Time) với sự tách rời tại các ranh giới chunk (Mục~\ref{sec:training}).

\subsection{Mã hóa Sự kiện và Điều biến Danh tính Tiến trình}
\label{sec:event_encoding}

Mỗi sự kiện kiểm toán $e_t$ mang ba loại đầu vào thô: loại sự kiện rời rạc (một trong $\sim$40 loại sự kiện CDM như \texttt{EVENT\_READ}, \texttt{EVENT\_MMAP}, \texttt{EVENT\_CONNECT}), một vectơ đặc trưng liên tục $\mathbf{x}_{\text{cont}} \in \mathbb{R}^{n_c}$ (số byte, cờ, thống kê thời gian đến), và đường dẫn hệ thống tệp của tiến trình chủ thể (ví dụ: \texttt{/usr/bin/bash}, \texttt{/home/admin/clean}). Bộ mã hóa sự kiện chiếu chúng vào một biểu diễn dày đặc $\mathbf{f}_e \in \mathbb{R}^{2d}$ dùng làm đầu vào cho tất cả các giai đoạn tiếp theo.

\textbf{Loại sự kiện và đặc trưng liên tục.} Loại sự kiện được ánh xạ tới một nhúng học được $\mathbf{e}_{\text{type}} = \text{Embed}_{\text{type}}(\text{type}(e_t)) \in \mathbb{R}^d$. Các đặc trưng liên tục được chiếu thông qua một lớp tuyến tính $\mathbf{c} = \mathbf{W}_c \mathbf{x}_{\text{cont}} + \mathbf{b}_c \in \mathbb{R}^d$.

\textbf{Điều biến danh tính tiến trình có cổng.} Các định danh thực thể số nguyên được gán trong quá trình tiền xử lý không có ý nghĩa ngữ nghĩa: cùng một tệp mã độc \texttt{/home/admin/clean} nhận các ID nguyên khác nhau trong các chiến dịch khác nhau vì từ vựng thực thể được xây dựng trên từng tập dữ liệu. Để cung cấp tín hiệu ngữ nghĩa khái quát hóa trên các chiến dịch, chúng tôi nhúng đường dẫn tiến trình của chủ thể thông qua phép tra cứu học được $\mathbf{p}_{\text{raw}} = \text{Embed}_{\text{proc}}(\text{path}(v_{\text{subj}})) \in \mathbb{R}^{64}$ từ một từ vựng về tên tiến trình đã biết (định nghĩa trong Mục~\ref{sec:setup}). Nó được chiếu tới chiều của mô hình $\mathbf{p} = \mathbf{W}_p \mathbf{p}_{\text{raw}} \in \mathbb{R}^d$.

Danh tính tiến trình được tiêm như một \emph{cổng thặng dư cộng} vào nhúng loại sự kiện thay vì nối lại:
\begin{equation}
    \mathbf{e}_{\text{mod}} = \mathbf{e}_{\text{type}} + g_p \cdot \mathbf{p}
    \label{eq:process_gate}
\end{equation}
trong đó $g_p = \sigma(\gamma_p)$ là một cổng vô hướng được kiểm soát bởi một tham số học duy nhất $\gamma_p$, khởi tạo bằng $-2.0$ sao cho $g_p \approx 0.12$ khi bắt đầu đào tạo. Danh tính tiến trình được tiêm vào nhúng loại sự kiện thay vì các đặc trưng liên tục vì đường dẫn tiến trình điều biến ngữ nghĩa sự kiện (ví dụ: một \texttt{EVENT\_MMAP} bởi \texttt{firefox} mang ý nghĩa ngữ cảnh khác với \texttt{pulseaudio}), trong khi các đặc trưng liên tục đo lường các đại lượng không phụ thuộc tiến trình như số byte. Khởi tạo cổng gần bằng 0 buộc mô hình phải học trước từ loại sự kiện và đặc trưng liên tục trước khi tín hiệu danh tính tiến trình được trộn vào. Nếu không có cổng, tín hiệu tiến trình cộng vào sẽ không bị ràng buộc tại thời điểm khởi tạo, tạo ra một lối tắt gradient qua danh tính tiến trình làm bỏ qua sự phân biệt cấp sự kiện. Chúng tôi xác nhận chế độ lỗi này theo thực nghiệm trong Mục~\ref{sec:ablation_process}.

\textbf{Dropout danh tính.} Trong quá trình đào tạo, danh tính tiến trình của mỗi sự kiện được thay thế độc lập bằng mã thông báo không xác định (chỉ mục 0) với xác suất $\rho = 0.3$:
\begin{equation}
    \text{path}(v_{\text{subj}}) \leftarrow 
    \begin{cases} 
        \text{path}(v_{\text{subj}}) & \text{với xác suất } 1 - \rho \\
        \texttt{\textless unk\textgreater} & \text{với xác suất } \rho
    \end{cases}
\end{equation}
Điều này ngăn bộ phân loại phụ thuộc vào danh tính tiến trình như một tín hiệu phát hiện đủ. Mô hình phải duy trì một con đường phát hiện khả thi chỉ thông qua các đặc trưng loại sự kiện và trạng thái thực thể, sử dụng danh tính tiến trình như một sự điều biến bổ sung thay vì một lối tắt.

\textbf{Biểu diễn sự kiện cuối cùng.} Ba thành phần được kết hợp và chuẩn hóa:
\begin{equation}
    \mathbf{f}_e = \text{LayerNorm}\!\left([\mathbf{e}_{\text{mod}} \;\|\; \mathbf{c}]\right) \in \mathbb{R}^{2d}
    \label{eq:event_repr}
\end{equation}
nơi $\|$ biểu thị phép nối. Loại sự kiện (được điều biến bởi danh tính tiến trình) chiếm $d$ chiều đầu tiên; các đặc trưng liên tục chiếm $d$ chiều thứ hai.

\subsection{Ngân hàng Trạng thái Thực thể}
\label{sec:entity_bank}

\texttt{EntityStateBank} lưu trữ hai tensor như các bộ đệm được đăng ký (không phải tham số học được): một ma trận trạng thái $\mathbf{S} \in \mathbb{R}^{|\mathcal{V}| \times d}$, trong đó hàng $\mathbf{S}[v]$ giữ trạng thái hiện tại $\mathbf{s}_v$ của thực thể $v$, và một vectơ dấu thời gian nhìn thấy lần cuối $\boldsymbol{\tau} \in \mathbb{R}^{|\mathcal{V}|}$, trong đó $\boldsymbol{\tau}[v]$ ghi lại thời gian gần đây nhất mà thực thể $v$ tham gia vào bất kỳ sự kiện nào. Cả hai tensor đều được khởi tạo về 0 (trạng thái) và $-1$ (dấu thời gian, biểu thị chưa thấy) vào đầu mỗi epoch đào tạo và duy trì qua các chunk sự kiện trong epoch. Đặt lại tại ranh giới epoch giúp mô hình không bị quá khớp vào quỹ đạo trạng thái cụ thể tích lũy ở các epoch trước, trong khi sự duy trì nội-epoch cung cấp bối cảnh tầm xa cần thiết bởi luận điểm.

\textbf{Giới hạn bộ nhớ.} Ngân hàng yêu cầu $|\mathcal{V}| \times (d + 1)$ giá trị dấu phẩy động. Đối với tập dữ liệu Theia ($|\mathcal{V}| \approx 7.1 \times 10^6$, $d = 128$), bộ nhớ này chiếm 3.6 GB ở dạng \texttt{float32}, nằm gọn trong một GPU duy nhất.

\textbf{Truy xuất trạng thái và thời gian trôi qua.} Đối với một sự kiện đến $e_t$ với các khe thực thể $(v_{\text{subj}}, v_{\text{obj}}, v_{\text{obj2}})$, mô hình thu thập một tensor trạng thái $\mathbf{S}_e \in \mathbb{R}^{3 \times d}$ bằng cách lập chỉ mục vào $\mathbf{S}$. Một mặt nạ hợp lệ $\mathbf{m} \in \{0, 1\}^3$ được đặt thành 0 cho các thực thể vắng mặt: các siêu cạnh kích thước 2 lưu $v_{\text{obj2}} = -1$, được giới hạn tại chỉ số 0 trong quá trình truy xuất. Mặt nạ đảm bảo rằng trạng thái được đọc sai tại chỉ số 0 sẽ đóng góp 0 vào cả tính toán tập hợp và đường ghi lại, ngăn sự phá hỏng trạng thái của thực thể 0.

Thời gian trôi qua cho mỗi thực thể được tính như sau:
\begin{equation}
    \Delta t_v = \max\!\left(0,\; t_{\text{curr}} - \boldsymbol{\tau}[v]\right)
    \label{eq:dt}
\end{equation}
nơi $\Delta t_v = 0$ cho các thực thể chưa thấy trước đó ($\boldsymbol{\tau}[v] = -1$). Giá trị tối thiểu ở mức 0 ngăn ngừa lỗi dấu thời gian không đơn điệu do sắp xếp lại bên trong shard. Tất cả các dấu thời gian được chuyển đổi sang giây so với tham chiếu toàn cục $t_0$ (dấu thời gian nano giây của sự kiện đầu tiên trong tập đào tạo):
\begin{equation}
    t_{\text{curr}} = (\text{timestamp\_nanos}(e_t) - t_0) \;/\; 10^9
    \label{eq:t_curr}
\end{equation}
Tất cả các bộ chia (đào tạo, xác thực, kiểm tra) chia sẻ cùng $t_0$ để $\Delta t_v$ vẫn hợp lệ qua ranh giới đào tạo-kiểm tra trong quá trình đánh giá khởi động ấm (warm-start) (Mục~\ref{sec:training}).

Thời gian trôi qua được chuyển đến các giai đoạn sau ở dạng tỷ lệ logarit là $\log(1 + \Delta t_v)$ để nén dải động: khoảng cách giữa các sự kiện trong tập dữ liệu Theia kéo dài từ micro giây đến hàng giờ.

\textbf{Chuẩn hóa thực thể trung tâm.} Các thực thể có bậc cao (ví dụ: \texttt{init} với PID 1, tham gia vào hơn $1.5 \times 10^5$ sự kiện trong tập dữ liệu Theia) tích lũy các vectơ trạng thái với chuẩn không giới hạn thông qua các bản cập nhật SSM lặp lại. Không có chuẩn hóa, các phép chiếu khóa trong giai đoạn tập hợp (Mục~\ref{sec:aggregation}) tạo ra điểm số chú ý gây tràn \texttt{float32} softmax ($\exp(x)$ cho $x > 88$). Để ngăn chặn điều này, các trạng thái truy xuất được chuẩn hóa thông qua chuẩn hóa lớp trước khi vào giai đoạn tập hợp. Việc chuẩn hóa này chỉ được áp dụng cho bản sao đã truy xuất $\hat{\mathbf{S}}_e$; ngân hàng $\mathbf{S}$ giữ lại các trạng thái chưa chuẩn hóa:
\begin{equation}
    \hat{\mathbf{S}}_e = \text{LayerNorm}(\mathbf{S}_e) \odot \mathbf{m}
    \label{eq:state_norm}
\end{equation}
Mặt nạ hợp lệ $\mathbf{m}$ được áp dụng sau khi chuẩn hóa để đảm bảo các khe thực thể vắng mặt đóng góp 0 vào các tính toán sau này. Tất cả các tham chiếu tiếp theo đến các trạng thái thực thể trong một lượt truyền thẳng đơn lẻ đều đề cập đến $\hat{\mathbf{S}}_e$.

\textbf{Ghi lại trạng thái.} Sau bước cập nhật trạng thái SSM (Mục~\ref{sec:ssm_update}), các trạng thái được cập nhật được phân tán trở lại vào $\mathbf{S}$ thông qua gán chỉ mục tại chỗ, và $\boldsymbol{\tau}$ được cập nhật với $t_{\text{curr}}$. Chỉ các thực thể hợp lệ (những thực thể có $\mathbf{m}_v = 1$) mới được ghi lại. Để ngăn hỏng hóc số học lan truyền qua các chunk, các giá trị NaN trong trạng thái được cập nhật được thay bằng 0 trước khi ghi lại.

\subsection{Tập hợp Siêu cạnh ($\mathcal{V} \to \mathcal{E}$)}
\label{sec:aggregation}

Cho trạng thái thực thể đã chuẩn hóa $\hat{\mathbf{S}}_e$ (Phương trình~\ref{eq:state_norm}) và biểu diễn sự kiện $\mathbf{f}_e$ (Phương trình~\ref{eq:event_repr}), giai đoạn tập hợp tạo ra một biểu diễn siêu cạnh duy nhất $\mathbf{x}_e \in \mathbb{R}^d$ kết hợp các trạng thái của 2--3 thực thể tham gia, được điều kiện hóa bởi ngữ nghĩa sự kiện. Việc tập hợp xác định \emph{mức độ} mà lịch sử tích lũy của mỗi thực thể đóng góp vào biểu diễn của sự kiện hiện tại, trực tiếp kiểm soát thông tin có sẵn cho bộ phân loại.

\textbf{Nhúng vai trò.} Mỗi khe thực thể được gán một nhúng vai trò học được $\mathbf{r}_v \in \mathbb{R}^d$ từ bảng tra cứu gồm 3 mục (Chủ thể, Đối tượng Chính, Đối tượng Phụ). Các nhúng vai trò là cần thiết bởi vì việc tập hợp hoạt động trên một tập hợp trạng thái không có thứ tự; nếu không có thông tin vai trò vị trí, cơ chế chú ý không thể phân biệt chủ thể với đối tượng, từ đó loại bỏ cấu trúc có hướng của sự kiện.

\textbf{Mô tả thực thể.} Đối với mỗi thực thể $v \in V_{e_t}$, chúng tôi xây dựng một bộ mô tả $\mathbf{h}_v$ bằng cách nối trạng thái chuẩn hóa, nhúng vai trò và thời gian trôi qua tỷ lệ logarit:
\begin{equation}
    \mathbf{h}_v = [\hat{\mathbf{s}}_v \;\|\; \mathbf{r}_v \;\|\; \log(1 + \Delta t_v)] \in \mathbb{R}^{2d + 1}
    \label{eq:entity_descriptor}
\end{equation}

\textbf{Chú ý động.} Các đặc trưng sự kiện $\mathbf{f}_e$ đóng vai trò là truy vấn, và các bộ mô tả thực thể đóng vai trò là khóa và giá trị. Điều này điều kiện hóa sự tập hợp vào loại sự kiện: đối với một \texttt{EVENT\_WRITE}, đối tượng (tệp đang được ghi) có thể mang nhiều trạng thái phân biệt hơn chủ thể, trong khi đối với một \texttt{EVENT\_MMAP}, đối tượng phụ (vùng bộ nhớ) mang tín hiệu được xác định trong Mục~\ref{sec:hypergraph}. Quá trình tính toán chú ý là:
\begin{align}
    \mathbf{q} &= \mathbf{W}_q \mathbf{f}_e \in \mathbb{R}^{d} \label{eq:agg_q} \\
    \mathbf{k}_v &= \mathbf{W}_k \mathbf{h}_v \in \mathbb{R}^{d} \label{eq:agg_k} \\
    \mathbf{v}_v &= \mathbf{W}_v \mathbf{h}_v \in \mathbb{R}^{d} \label{eq:agg_v}
\end{align}
trong đó $\mathbf{W}_q \in \mathbb{R}^{d \times 2d}$, $\mathbf{W}_k, \mathbf{W}_v \in \mathbb{R}^{d \times (2d+1)}$. Trọng số chú ý được tính qua tích vô hướng có tỷ lệ và mặt nạ. Đối với mỗi khe thực thể $v$, điểm số thô là:
\begin{equation}
    s_v = \mathbf{q}^\top \mathbf{k}_v \;/\; \sqrt{d}
    \label{eq:raw_score}
\end{equation}
Các khe không hợp lệ ($\mathbf{m}_v = 0$) có điểm được thiết lập thành $-\infty$ trước softmax, do đó $\exp(-\infty) = 0$ ở cả tử số và mẫu số:
\begin{equation}
    \alpha_v = 
    \begin{cases}
        \displaystyle\frac{\exp(s_v)}{\sum_{u:\,\mathbf{m}_u = 1} \exp(s_u)} & \text{nếu } \mathbf{m}_v = 1 \\[6pt]
        0 & \text{nếu } \mathbf{m}_v = 0
    \end{cases}
    \label{eq:attn_weights}
\end{equation}
Như một biện pháp bảo vệ, bất kỳ giá trị NaN nào trong các trọng số chú ý (có thể phát sinh từ các sự kiện suy biến khi tất cả các khe bị che lấp, tạo ra $0/0$ trong softmax) đều được thay thế bằng 0. Biểu diễn siêu cạnh lúc đó là:
\begin{equation}
    \mathbf{x}_e = \mathbf{W}_{\text{out}} \sum_{v \in V_{e_t}} \alpha_v \, \mathbf{v}_v \in \mathbb{R}^d
    \label{eq:x_e}
\end{equation}
với $\mathbf{W}_{\text{out}} \in \mathbb{R}^{d \times d}$ là phép chiếu đầu ra tuyến tính. Nhúng vai trò $\mathbf{r}_v$ được giữ lại và chuyển qua giai đoạn cập nhật SSM (Mục~\ref{sec:ssm_update}), tại đây chúng đảm bảo sự truyền bá trạng thái đặc thù cho từng vai trò.

\subsection{Cập nhật Trạng thái SSM Chọn lọc ($\mathcal{E} \to \mathcal{V}$)}
\label{sec:ssm_update}

Được cung cấp biểu diễn siêu cạnh $\mathbf{x}_e$ (Phương trình~\ref{eq:x_e}) và nhúng vai trò $\mathbf{r}_v$ từ giai đoạn tập hợp, bộ cập nhật SSM tính toán trạng thái mới $\mathbf{s}_v^{+}$ cho từng thực thể tham gia $v \in V_{e_t}$. Đây là cơ chế mà thông qua đó bối cảnh hành vi tích lũy truyền qua các thực thể và theo thời gian: trạng thái được cập nhật mã hóa cả lịch sử trước đó của thực thể (thông qua phần suy giảm) và đóng góp của sự kiện hiện tại (thông qua phần đầu vào).

\textbf{Ma trận suy giảm trạng thái.} Ma trận suy giảm trạng thái đường chéo $A \in \mathbb{R}^d$ được lưu trữ trong không gian log như một tham số có thể học được $\mathbf{A}_{\log} \in \mathbb{R}^d$ và được phục hồi như sau:
\begin{equation}
    A = -\exp\!\left(\text{clamp}(\mathbf{A}_{\log},\, -4,\, 4)\right)
    \label{eq:A}
\end{equation}
Tất cả các thành phần của $A$ đều âm nghiêm ngặt, đảm bảo các trạng thái phân rã theo thời gian khi không có đầu vào mới. $\mathbf{A}_{\log}$ được khởi tạo bằng $\log(\text{linspace}(0.1, 2.0, d))$, tạo ra tỷ lệ suy giảm ban đầu kéo dài $A_j \in [-2.0, -0.1]$. Sự khởi tạo đa quy mô thời gian này gán cho các chiều trạng thái khác nhau các tỷ lệ suy giảm khác nhau: các chiều phân rã chậm ($A_j \approx -0.1$) giữ thông tin qua nhiều sự kiện cho phụ thuộc tầm xa, trong khi các chiều phân rã nhanh ($A_j \approx -2.0$) đáp ứng với hoạt động gần đây. Việc giới hạn $\mathbf{A}_{\log}$ ở $[-4, 4]$ ngăn tràn số (overflow) trong phép lấy mũ.

\textbf{Phép chiếu đầu vào đặc thù theo vai trò.} Ma trận đầu vào $B$ được tính từ sự nối của biểu diễn siêu cạnh $\mathbf{x}_e$ và nhúng vai trò của thực thể $\mathbf{r}_v$:
\begin{equation}
    B_v = \mathbf{W}_B [\mathbf{x}_e \;\|\; \mathbf{r}_v] \in \mathbb{R}^d
    \label{eq:B}
\end{equation}
trong đó $\mathbf{W}_B \in \mathbb{R}^{d \times 2d}$. Việc bao gồm nhúng vai trò tạo ra các giá trị $B_v$ khác nhau cho chủ thể, đối tượng chính, và đối tượng phụ của cùng một sự kiện. Nếu không có điều này, cả ba thực thể sẽ nhận các hướng cập nhật giống hệt nhau từ $\mathbf{x}_e$, khiến các trạng thái của chúng hội tụ và phá hủy cấu trúc đồ thị có hướng phân biệt chủ thể và đối tượng.

\textbf{Rời rạc hóa nhận biết thời gian.} Các tham số SSM thời gian liên tục $(A, B_v)$ được rời rạc hóa sử dụng kích thước bước phụ thuộc vào đầu vào $\Delta_v \in \mathbb{R}^d$:
\begin{equation}
    \Delta_v = \text{clamp}\!\left(\text{softplus}\!\left(\mathbf{W}_\Delta [\mathbf{x}_e \;\|\; \mathbf{r}_v]\right) \cdot \left(1 + \lambda \log(1 + \Delta t_v)\right),\; \max\!=5\right)
    \label{eq:delta}
\end{equation}
nơi $\mathbf{W}_\Delta \in \mathbb{R}^{d \times 2d}$ với độ lệch (bias) khởi tạo ở $-2.0$ (tạo ra kích thước bước ban đầu là $\text{softplus}(-2) \approx 0.13$), $\lambda = 0.1$, và $\Delta t_v$ là thời gian trôi qua từ Mục~\ref{sec:entity_bank}. Phép nhân chia tỷ lệ thời gian $(1 + 0.1 \cdot \log(1 + \Delta t_v))$ áp dụng một sự điều chỉnh nhẹ (khoảng $1.0$ đến $1.7\times$ trên phạm vi quan sát được) giúp mở rộng bước rời rạc hóa cho các thực thể đã không hoạt động, cho phép sự suy giảm phản ánh thời gian thực đã trôi qua mà không gây ra sự mất ổn định gradient như việc nhân trực tiếp $\Delta t_v$ có thể gây ra. Giới hạn tại 5 ngăn chặn sự tràn số học trong bước lũy thừa tiếp theo.

Các tham số rời rạc hóa là:
\begin{align}
    \bar{A}_v &= \exp(\Delta_v \odot A) \in (0, 1)^d \label{eq:A_bar} \\
    \bar{B}_v &= \Delta_v \odot B_v \label{eq:B_bar}
\end{align}
trong đó $\bar{A}_v$ kiểm soát sự giữ lại trạng thái trước đó trên mỗi chiều (giá trị gần 1 là giữ lại, gần 0 là quên đi), và $\bar{B}_v$ tỷ lệ hóa mức độ đóng góp của đầu vào. Phép cập nhật SSM thô là:
\begin{equation}
    \tilde{\mathbf{s}}_v = \bar{A}_v \odot \hat{\mathbf{s}}_v + \bar{B}_v \odot \mathbf{x}_e
    \label{eq:ssm_raw}
\end{equation}

\textbf{Kết nối thặng dư có cổng.} Một cổng học được $\mathbf{g}_v \in (0, 1)^d$ trộn lẫn kết quả của SSM với trạng thái trước đó:
\begin{align}
    \mathbf{g}_v &= \sigma\!\left(\mathbf{W}_g [\mathbf{x}_e \;\|\; \mathbf{r}_v]\right) \label{eq:gate} \\
    \mathbf{s}_v^{+} &= \mathbf{g}_v \odot \tilde{\mathbf{s}}_v + (1 - \mathbf{g}_v) \odot \hat{\mathbf{s}}_v \label{eq:gated_update}
\end{align}
trong đó $\mathbf{W}_g \in \mathbb{R}^{d \times 2d}$ với bias được khởi tạo là $-3.0$, do đó $\mathbf{g}_v \approx 0.04$ lúc bắt đầu đào tạo. Việc khởi tạo gần đóng này là cần thiết về mặt cấu trúc: ở lần khởi tạo ngẫu nhiên, các tham số SSM $(A, B)$ sinh ra các bản cập nhật trạng thái tùy ý, và một cổng mở sẽ ngay lập tức ghi đè lên các trạng thái thực thể bằng nhiễu (noise). Khởi tạo cổng gần bằng không buộc mô hình chủ động học \emph{khi nào} nên truyền bá trạng thái, ngăn sự quá mượt (over-smoothing) ở bước khởi tạo. Cổng mở dần trong quá trình đào tạo khi các tham số SSM trở nên có ý nghĩa.

\subsection{Đầu Phân loại}
\label{sec:classification}

Giai đoạn cuối cùng quyết định xem sự kiện $e_t$ có phải là một phần của chuỗi tấn công hay không bằng cách đánh giá cả ngữ nghĩa độc lập của nó và trạng thái hành vi tích lũy của các thực thể tham gia.

\textbf{Định tuyến Gradient qua các trạng thái sau cập nhật.} Bộ phân loại hoạt động trên trạng thái \emph{sau cập nhật} $\mathbf{s}_v^{+}$ tính ở Mục~\ref{sec:ssm_update}, không phải trạng thái trước cập nhật được lấy từ ngân hàng. Định tuyến này cực kỳ quan trọng đối với cấu trúc cho việc đào tạo toàn trình (end-to-end). Các trạng thái của ngân hàng được coi là các hằng số tách rời trong quá trình truyền thẳng; nếu bộ phân loại đọc trực tiếp từ ngân hàng, biểu đồ tính toán sẽ kết thúc, cắt đứt hoàn toàn dòng gradient tới bộ cập nhật SSM (thứ tạo ra trạng thái mới) và tới bộ tập hợp siêu cạnh (thứ tạo ra đầu vào SSM). Bằng cách phân loại trên trạng thái sau cập nhật, gradient được truyền ngược qua toàn bộ chuỗi truyền thẳng.

\textbf{Lấy trung bình có mặt nạ hợp lệ.} Trạng thái sau cập nhật được gộp (pooled) vào một vectơ bối cảnh mức độ sự kiện duy nhất $\bar{\mathbf{s}} \in \mathbb{R}^d$. Quá trình này yêu cầu che mờ các đối tượng phụ vắng mặt (nơi $\mathbf{m}_v = 0$):
\begin{equation}
    \bar{\mathbf{s}} = \text{LayerNorm}\!\left( \frac{1}{\sum_{v} \mathbf{m}_v} \sum_{v \in V_{e_t}} \mathbf{m}_v \cdot \mathbf{s}_v^{+} \right)
\end{equation}
Việc che mờ trước khi cộng là bắt buộc do cập nhật SSM (Phương trình~\ref{eq:ssm_raw}) sinh ra bản cập nhật khác 0 ngay cả đối với khe thực thể vô giá trị (được điều khiển bởi biểu diễn sự kiện $\mathbf{x}_e$ thông qua phép chiếu vai trò $\bar{B}_v$). Không có mặt nạ, khe rỗng sẽ làm hỏng giá trị trung bình. Trạng thái gộp được chuẩn hóa lớp để làm ổn định đầu vào bộ phân loại.

\textbf{Xác suất phát hiện.} Trạng thái đã gộp được nối với các đặc trưng sự kiện đã chiếu $\mathbf{W}_p \mathbf{f}_e$ (với $\mathbf{W}_p \in \mathbb{R}^{d \times 2d}$). Xác suất $\hat{y}_t \in [0, 1]$ cho sự kiện $e_t$ là sự kiện tấn công được tính thông qua Multi-Layer Perceptron (MLP):
\begin{equation}
    \hat{y}_t = \sigma\!\left( \text{MLP}\!\left(\left[\bar{\mathbf{s}} \;\|\; \mathbf{W}_p \mathbf{f}_e\right]\right) \right)
\end{equation}
MLP bao gồm ba lớp tuyến tính với hàm kích hoạt GELU trung gian và Dropout ($\rho=0.1$) sau lớp đầu tiên, chiếu số chiều từ $2d \to d \to d/2 \to 1$. Hàm mất mát Binary Cross-Entropy (BCE) được tính toán theo nhãn thực tế của sự kiện $y_t \in \{0, 1\}$.

\subsection{Đào tạo: TBPTT và Khởi động ấm Liên tục}
\label{sec:training}

HyperMamba yêu cầu các kỹ thuật đào tạo khác với mô hình NIDS không trạng thái vì các trạng thái thực thể tồn tại dai dẳng qua chuỗi sự kiện.

\textbf{Lan truyền ngược qua Thời gian Bị cắt ngắn (TBPTT).} Để xử lý các luồng sự kiện liên tục mà không làm cạn kiệt bộ nhớ, các sự kiện được phân tách theo thời gian thành các đoạn (chunk) kích thước $C = 4096$. Với mỗi chunk, mô hình thực hiện truyền thẳng (Mục~\ref{sec:event_encoding}--\ref{sec:classification}) và tính mất mát BCE. Quan trọng là, các trạng thái cập nhật được ghi vào \texttt{EntityStateBank} được giữ lại cho chunk tiếp theo, nhưng biểu đồ tính toán của chúng bị tách ra qua lệnh `.detach()`. Chiến lược TBPTT này cho phép bối cảnh \emph{truyền tới} tồn tại vô hạn trên toàn bộ tập dữ liệu, tích lũy các tín hiệu nhiễu tầm xa, trong khi giới hạn tính toán truyền ngược gradient ở mức $\mathcal{O}(C)$.

\textbf{Tối ưu hóa và sự ổn định.} Mô hình được tối ưu bằng Adam với lịch trình tỷ lệ học Cosine Annealing. Để xử lý sự mất cân bằng lớp nghiêm trọng trong luồng kiểm toán, tổn thất BCE được tính theo tỷ trọng số lượng mẫu âm so với mẫu dương trong tập đào tạo (giới hạn ở mức $30.0$ để tránh bùng nổ gradient). Vì SSM hoạt động trên các chuỗi rời rạc đôi khi có thể tạo ra tính bất ổn số học, các chunk dẫn đến tổn thất NaN hoặc gradient NaN sẽ bị bỏ qua và ngân hàng trạng thái ngay lập tức bị tách ra (detach) để ngăn chặn làm hỏng các chunk sau. Chuẩn Gradient cũng được cắt gọn ở mức 1.0.

\textbf{Cài đặt lại ranh giới Epoch.} \texttt{EntityStateBank} hoàn toàn được đưa về 0 ở đầu mỗi epoch đào tạo. Nếu không thiết lập lại, mô hình sẽ bắt đầu epoch $k$ (đào tạo trên phân đoạn thời gian đầu tiên) sử dụng các trạng thái thực thể tương lai tích lũy ở cuối epoch $k-1$ (sau khi đánh giá trên các shard xác thực). Điều này sẽ tạo ra rò rỉ dữ liệu thời gian nghiêm trọng. Thiết lập lại đảm bảo mỗi epoch học cách tích lũy trạng thái hoàn toàn từ quá khứ.

\textbf{Đánh giá khởi động ấm liên tục.} Trong một triển khai thực tế, NIDS hoạt động liên tục; trạng thái thực thể không bao giờ "lạnh" (trống). Để đánh giá chính xác khả năng phát hiện trên một tập dữ liệu kiểm tra, chúng ta phải mô phỏng triển khai liên tục này. Trước khi đánh giá tập kiểm tra, mô hình chạy quá trình \emph{khởi động (warmup)} chỉ truyền thẳng (tắt gradient và việc thu thập số liệu) đi qua tất cả phân đoạn đào tạo và xác thực trước đó theo thứ tự thời gian nghiêm ngặt để lấp đầy \texttt{EntityStateBank}. Sau đó, tập kiểm tra được đánh giá sử dụng các trạng thái đã "ấm" này. Việc đánh giá một mô hình có trạng thái từ khởi động lạnh (cold start) sẽ đánh giá thấp một cách hệ thống khả năng phát hiện các cuộc tấn công nhiều giai đoạn bắt đầu từ trước cửa sổ kiểm tra.

\textbf{Lựa chọn ngưỡng.} Theo tiêu chuẩn đánh giá nghiêm ngặt cho học máy bảo mật~\cite{arp2022dos}, ngưỡng ra quyết định được chọn động để tối đa hóa F1-score \emph{độc quyền} trên tập xác thực. Ngưỡng cố định này sau đó được áp dụng cho tập kiểm tra, ngăn ngừa hoàn toàn sự rò rỉ tối ưu hóa ngưỡng của tập kiểm tra.

\section{Đánh giá}
\label{sec:evaluation}

\subsection{Thiết lập Thử nghiệm}
\label{sec:experimental_setup}

\textbf{Tập dữ liệu và định nghĩa nhãn.} Chúng tôi đánh giá HyperMamba trên tập dữ liệu DARPA Transparent Computing (TC) Engagement 3 Theia, một thước đo tiêu chuẩn cho NIDS dựa trên nguồn gốc. Sau khi xử lý, tập dữ liệu bao gồm 25 phân đoạn theo trình tự thời gian (shards) chứa khoảng 105 triệu sự kiện. Chúng tôi nhắm mục tiêu vào nhãn \emph{crossprocess}: các sự kiện do các tiến trình con nhiều giai đoạn thực thi được sinh ra từ tiến trình đầu vào ban đầu bị xâm phạm. Bản thân tiến trình điểm vào không nằm trong nhãn này, buộc hệ thống phát hiện phải định danh hoạt động độc hại ở các giai đoạn sau (ví dụ: tập lệnh kết nối đến máy chủ C2) thay vì chỉ đánh dấu hành vi dị thường ban đầu của ứng dụng bị xâm phạm.

\textbf{Phân tách (Splits).} Việc đánh giá hoàn toàn theo trình tự thời gian để tránh rò rỉ dữ liệu. Dữ liệu bao gồm hai chiến dịch APT đa giai đoạn riêng biệt. Chiến dịch 1 kéo dài từ shard 0--10 và Chiến dịch 2 kéo dài từ shard 12--24. Để đánh giá sự tiếp nối trong cùng chiến dịch và sự khái quát hóa giữa các chiến dịch, chúng tôi xác định hai chế độ đánh giá:
\begin{itemize}
    \item \textbf{Chiến dịch hỗn hợp (\texttt{full}):} Cả hai chiến dịch đều có mặt trong tập đào tạo. Chúng tôi gán các shard $\{0..7, 12..19\}$ cho đào tạo, các shard $\{8, 9, 20, 21\}$ cho xác thực, và các shard $\{10, 22..24\}$ cho kiểm tra. Điều này khớp với quy trình đánh giá chuẩn của các phương pháp cơ sở.
    \item \textbf{Chéo chiến dịch (\texttt{cross}):} Kiểm tra khả năng khái quát cho các chiến lược tấn công chưa từng thấy. Mô hình được huấn luyện hoàn toàn trên Chiến dịch 1 (đào tạo shards $\{0..8\}$, xác thực shards $\{9, 10\}$) và kiểm tra hoàn toàn trên Chiến dịch 2 (shards kiểm tra $\{12..24\}$).
\end{itemize}

\textbf{Mô hình cơ sở (Baselines).} Chúng tôi so sánh HyperMamba với NIDS tiên tiến dựa trên nguồn gốc, định vị mỗi hệ thống theo độ phân giải phát hiện cốt lõi của chúng:
\begin{itemize}
    \item \textbf{KAIROS}~\cite{kairos}: Hoạt động ở cấp độ đồ thị, gán điểm dị thường cho toàn bộ các đồ thị nguồn gốc khổng lồ.
    \item \textbf{MAGIC}~\cite{magic}: Hoạt động ở cấp độ đồ thị, phát hiện tấn công thông qua lỗi tái cấu trúc của các đồ thị con cấu trúc.
    \item \textbf{FLASH}~\cite{flash}: Hoạt động ở cấp độ nút (node), cảnh báo tiến trình hoặc tệp dị thường.
    \item \textbf{ThreaTrace}~\cite{threatrace}: Hoạt động ở cấp độ nút, phân loại các thực thể bên trong đồ thị nguồn gốc.
\end{itemize}
Điểm then chốt là không mô hình cơ sở nào có khả năng phát hiện \emph{cấp độ sự kiện} nguyên bản, vốn là mức chi tiết nhỏ nhất và là đầu ra chính của HyperMamba. 

\textbf{Số liệu và Phương pháp ngưỡng.} Chúng tôi báo cáo Area Under the Precision-Recall Curve (AUPRC) như số liệu chính không phụ thuộc ngưỡng, do sự mất cân bằng dữ liệu nghiêm trọng. Đối với số liệu tính theo ngưỡng (F1-score và Tỷ lệ Dương tính Giả - FPR), chúng tôi tuân thủ hướng dẫn khắt khe dành cho AI bảo mật bởi Arp và cộng sự~\cite{arp2022dos}: ngưỡng được chọn động từ dự đoán xác thực dùng điểm vận hành tối đa F1, sau đó được cố định áp dụng cho tập dữ liệu kiểm tra. Điều này ngăn chặn nghiêm ngặt sự rò rỉ tối ưu hóa ngưỡng kiểm tra. Số liệu được báo cáo ở cả cấp độ sự kiện (với HyperMamba) và cấp độ nút (với HyperMamba và baselines).

\subsection{Kết quả Chính (C4)}
\label{sec:main_results}

Chúng tôi báo cáo kết quả phát hiện của HyperMamba so với các baselines trong Bảng~\ref{tab:main_results}.

\textbf{Sự không thể so sánh về độ phân giải.} Chúng tôi nhấn mạnh rằng việc so sánh bằng số trực tiếp giữa số liệu cấp sự kiện của HyperMamba và số liệu cấp nút hoặc đồ thị của các phương pháp khác là không hợp lệ về mặt cấu trúc. Phát hiện một nút bất thường trong $10^4$ tiến trình là một bài toán thống kê cơ bản khác và dễ hơn nhiều so với việc phân loại chính xác $10$ sự kiện tấn công trong số $10^5$ sự kiện mà nút đó thực thi. Chúng tôi báo cáo số liệu cấp sự kiện và cấp nút được tập hợp để chứng minh rằng HyperMamba vẫn cạnh tranh tốt ở cấp độ nút thô trong khi cung cấp độc quyền độ chính xác pháp y cấp sự kiện chi tiết hơn.

\begin{table*}[t]
    \centering
    \caption{Hiệu suất Phát hiện trên DARPA TC E3 Theia. HyperMamba là hệ thống duy nhất có khả năng phát hiện cấp sự kiện. Các con số cơ sở được báo cáo ở các mức phân giải gốc, thô hơn của chúng.}
    \label{tab:main_results}
    \begin{tabular}{l|c|ccc|cc}
        \hline
        \textbf{Model} & \textbf{Độ phân giải} & \textbf{Event AUPRC} & \textbf{Event F1} & \textbf{Event FPR} & \textbf{Node AUPRC} & \textbf{Node F1} \\
        \hline
        KAIROS & Graph & - & - & - & - & - \\
        MAGIC & Graph & - & - & - & - & - \\
        FLASH & Node & - & - & - & - & - \\
        ThreaTrace & Node & - & - & - & - & - \\
        \hline
        HyperMamba (Cùng chiến dịch) & Event \& Node & 0.86 & - & 0.000010 & - & - \\
        HyperMamba (Chéo chiến dịch) & Event \& Node & 0.72 & - & 0.000195 & - & - \\
        \hline
    \end{tabular}
\end{table*}

\textbf{Phát hiện Cùng chiến dịch so với Chéo chiến dịch.} Chúng tôi tách hiệu suất của HyperMamba thành hai chế độ. Trong điều kiện tiếp tục \emph{cùng chiến dịch} (same-campaign) (đánh giá ở các giai đoạn sau chưa thấy của chiến dịch có trong tập đào tạo), HyperMamba đạt AUPRC cấp sự kiện là $0.86$ với tỷ lệ Dương tính giả gần 0 là $0.000010$. Ở quy mô 100 triệu sự kiện, con số này chỉ tạo ra khoảng $1,000$ cảnh báo sai --- một khối lượng có thể quản lý được đối với các nhà phân tích, tương phản sâu sắc với sự mệt mỏi do cảnh báo từ các hệ thống có FPR cao.

Trong cài đặt tổng quát hóa \emph{chéo chiến dịch} khắt khe (chỉ đào tạo trên Chiến dịch 1 và kiểm tra trên Chiến dịch 2), AUPRC cấp sự kiện vẫn mạnh ở mức $0.72$, với FPR là $0.000195$. Điều này chứng minh rằng cơ chế truyền trạng thái SSM học được các hành vi tấn công có thể khái quát hóa thay vì chỉ đơn giản là ghi nhớ cấu trúc mạng của chiến dịch đào tạo.

\subsection{Cắt bỏ: Truyền Trạng thái Thực thể (C2)}
\label{sec:ablation_state}

Luận điểm trung tâm của HyperMamba là trạng thái thực thể tích lũy liên tục --- không phải các đặc điểm sự kiện riêng biệt --- là cơ chế chính cho phép phát hiện các cuộc tấn công tàng hình xuyên tiến trình. Để chứng minh theo thực nghiệm (Đóng góp 2), chúng tôi so sánh cấu trúc đầy đủ của HyperMamba với phiên bản cắt bỏ (ablation) không có trạng thái (\texttt{no\_state}).

\textbf{Thiết kế cắt bỏ.} Biến thể \texttt{no\_state} giữ nguyên kiến trúc: bộ mã hóa sự kiện, bộ tập hợp siêu cạnh, và đầu phân loại không đổi. Thay đổi duy nhất là ngăn cấu trúc tích lũy trạng thái tầm xa bằng cách thiết lập lại \texttt{EntityStateBank} ở mỗi ranh giới chunk, thay vì tách rời và duy trì các trạng thái. Mô hình buộc phải đánh giá mỗi chunk chỉ dùng các đặc trưng sự kiện cục bộ và bối cảnh bên trong chunk.

\begin{table}[t]
    \centering
    \caption{Cắt bỏ: Tác động của việc Truyền Trạng thái Thực thể Liên tục}
    \label{tab:ablation_state}
    \begin{tabular}{l|cc|cc}
        \hline
        \textbf{Model} & \textbf{Train AUPRC} & \textbf{Train F1} & \textbf{Test AUPRC} & \textbf{Test F1} \\
        \hline
        HyperMamba (Full) & 0.86 & - & 0.73 & - \\
        HyperMamba (\texttt{no\_state}) & 0.51 & - & 0.33 & - \\
        \hline
    \end{tabular}
\end{table}

\textbf{Kết quả.} Bảng~\ref{tab:ablation_state} cho thấy, loại bỏ truyền trạng thái dẫn đến sụp đổ thảm khốc về khả năng phát hiện. AUPRC kiểm tra giảm từ $0.73$ xuống $0.33$. Đáng chú ý, AUPRC huấn luyện cũng sụp đổ từ $0.86$ xuống $0.51$. Sự thất bại trong việc học ngay cả với phân phối huấn luyện xác nhận lập luận cấu trúc ở Mục~\ref{sec:threat_model}: sự kiện tấn công xuyên tiến trình không thể nhận dạng được nếu chỉ có đặc điểm sự kiện cục bộ. Nếu không có lịch sử tích lũy của tiến trình xâm phạm ban đầu, các tiến trình con không thể phân biệt được bằng thống kê so với tiến trình bình thường. Điều này chứng nhận luận điểm của HyperMamba: các cách tiếp cận không trạng thái trước đây không thể đối sánh được với độ chuẩn xác này.

\begin{figure}[t]
    \centering
    % \includegraphics[width=\linewidth]{figures/training_curves.pdf}
    \caption{Đường cong AUPRC huấn luyện và xác thực theo epoch. Phiên bản \texttt{no\_state} phân kỳ ngay lập tức, thất bại trong việc tìm thấy tín hiệu phân biệt.}
    \label{fig:training_curves}
\end{figure}

Biểu đồ (Hình~\ref{fig:training_curves}) thể hiện thêm cơ chế này. AUPRC xác thực của bản đầy đủ cải thiện đều đặn. Ngược lại, phiên bản \texttt{no\_state} thất bại hoàn toàn ngay từ đầu trong việc vạch ranh giới phân loại.

\subsection{Cắt bỏ: Tập hợp Chéo Thực thể}
\label{sec:ablation_cross_entity}

Trong khi phần trước chứng minh trạng thái là bắt buộc, ta cũng cần đánh giá xem việc \emph{chia sẻ} trạng thái này qua các thực thể trong một sự kiện có cần thiết không. Để xác định xem cấu trúc siêu cạnh (Mục~\ref{sec:aggregation}) có ích gì ngoài việc theo dõi riêng lẻ (Đóng góp 1), chúng tôi đánh giá \texttt{no\_cross\_entity}.

\textbf{Thiết kế cắt bỏ.} Biến thể này tắt tập hợp siêu cạnh. SSM chỉ cập nhật chủ thể dùng trạng thái cũ của nó. Trạng thái đối tượng không ảnh hưởng đến, và không nhận được vết tích lũy của chủ thể.

\begin{table}[h]
    \centering
    \caption{Cắt bỏ: Tác động của Tập hợp Siêu cạnh Chéo Thực thể}
    \label{tab:ablation_cross_entity}
    \begin{tabular}{l|cc}
        \hline
        \textbf{Model} & \textbf{Event AUPRC} & \textbf{Event F1} \\
        \hline
        HyperMamba (Full) & 0.86 & - \\
        HyperMamba (\texttt{no\_cross\_entity}) & - & - \\
        \hline
    \end{tabular}
\end{table}

\textbf{Kết quả.} Theo Bảng~\ref{tab:ablation_cross_entity}, việc tắt tập hợp làm giảm độ phát hiện. Dù tốt hơn không có trạng thái, nó không thể đạt được kết quả như mô hình siêu đồ thị toàn diện. Điều này xác nhận rằng theo dõi luồng topo qua các hoạt động mang lại tín hiệu thiết yếu để tìm ra các chiến dịch tinh vi.

\subsection{Cắt bỏ: Điều biến Danh tính Tiến trình (C3)}
\label{sec:ablation_process}

Dù sự truyền bá trạng thái cho phép mô hình liên kết sự kiện, việc khái quát nó sang chiến dịch mới yêu cầu các điểm neo ngữ nghĩa. Do các ID nguyên không có ý nghĩa qua các chiến dịch (ID $e_1$ có thể là `bash` ở đây, nhưng là tệp nhị phân độc hại ở kia), để chứng minh rằng danh tính tiến trình có cổng là yếu tố then chốt cho khái quát hóa (C3), chúng tôi cắt bỏ nhúng đường dẫn trên đánh giá \emph{chéo chiến dịch}.

\begin{table}[h]
    \centering
    \caption{Cắt bỏ: Điều biến Danh tính Tiến trình trên Tập Chéo Chiến dịch}
    \label{tab:ablation_process}
    \begin{tabular}{l|c}
        \hline
        \textbf{Model} & \textbf{Node AUPRC} \\
        \hline
        HyperMamba (Full) & 0.89 \\
        HyperMamba (\texttt{no\_process\_identity}) & - \\
        \hline
    \end{tabular}
\end{table}

\textbf{Kết quả.} Bảng~\ref{tab:ablation_process} cho thấy hiệu ứng. Không có nhúng tiến trình điều biến bộ mã hóa, mô hình hoàn toàn phải dựa vào mô hình thời gian. Điều này làm suy giảm nghiêm trọng khả năng tổng quát. Các nhúng định danh cung cấp một tín hiệu bất biến cốt lõi, giúp mô hình nhận biết được cấu trúc bất thường của `bash` khác cơ bản với mô hình đó của `nginx`.

\subsection{Chi phí Tính toán}
\label{sec:computational_cost}

Một NIDS phải có khả năng thời gian thực. Chúng tôi đánh giá hiệu suất xử lý của HyperMamba ở Bảng~\ref{tab:comp_cost}.

\begin{table}[h]
    \centering
    \caption{Chi phí Tính toán và Thông lượng}
    \label{tab:comp_cost}
    \begin{tabular}{l|c}
        \hline
        \textbf{Metric} & \textbf{Value} \\
        \hline
        Tham số (HyperMamba) & 1.17M \\
        Tham số (GRU-TGN Baseline) & 1.31M \\
        \hline
        Bộ nhớ Trạng thái ($7.1\text{M} \times 256$) & 7.2 GB \\
        \hline
        Thông lượng Train (GPU) & 734K sự kiện/giây \\
        Thông lượng Suy luận (GPU) & 1.9M sự kiện/giây \\
        Thông lượng Suy luận (CPU) & 70K sự kiện/giây \\
        \hline
    \end{tabular}
\end{table}

\textbf{Thông lượng.} Với tốc độ đỉnh $10,000$ sự kiện/giây trên một máy Theia, HyperMamba đạt thông lượng $1.94$M sự kiện/giây (GPU), gấp $194\times$ mức tải. Thậm chí chỉ với CPU, HyperMamba xử lý $70$K sự kiện/giây, vẫn gấp $7\times$ mức tải.

\textbf{Bộ nhớ.} Mô hình rất tối ưu số lượng tham số (1.17M). \texttt{EntityStateBank} cần $7.2$ GB bộ nhớ cho $7.1$ triệu thực thể, vừa vặn hoàn toàn với GPU phổ thông.

\section{Thảo luận}
\label{sec:discussion}

Chúng tôi nêu các giới hạn của nghiên cứu này như là những ràng buộc khách quan.

\textbf{Phạm vi Dữ liệu.} Chỉ thử nghiệm trên DARPA TC E3. Sự tổng quát đối với Linux/FreeBSD khác hoặc eBPF chưa được xác thực.

\textbf{Phạm vi Nhãn.} Tấn công không sinh ra tiến trình (e.g. lỗ hổng cấu trúc bộ nhớ nội bộ) nằm ngoài nhãn crossprocess và hiện tại chưa được thử nghiệm.

\textbf{Sự rời đi của Thực thể và LRU.} Cấu trúc bộ nhớ hỗ trợ xóa LRU nhưng chưa thử nghiệm căng thẳng (stress-test). Tác động của LRU lên kết quả vẫn là một câu hỏi mở.

\textbf{Tính bền vững Adversarial.} Giả định là black-box. Tấn công white-box pha loãng trạng thái (inject sự kiện thường để rửa vết độc) chưa được xem xét.

\textbf{Tấn công dựa trên Tải trọng (Payload).} Hệ thống hoạt động trên siêu dữ liệu kiểm toán (audit metadata). Tấn công payload, chẳng hạn như SQL injection, nằm ngoài cấu trúc.

\textbf{Phạm vi Máy chủ Đơn (Single-host).} Chưa xử lý sự di chuyển ngang (lateral movement) qua các máy chủ trên toàn mạng lưới lớn.

\section{Nghiên cứu Liên quan}
\label{sec:related_work}

Bảng~\ref{tab:positioning} tóm tắt các điểm tương quan.

\begin{table}[t]
    \centering
    \caption{Định vị HyperMamba so với NIDS dựa trên nguồn gốc qua 3 trục.}
    \label{tab:positioning}
    \begin{tabular}{l|ccc}
        \hline
        \textbf{Hệ thống} & \textbf{Độ phân giải} & \textbf{Trạng thái} & \textbf{Giám sát} \\
        \hline
        KAIROS & Graph & Snapshot & Tự Giám sát \\
        MAGIC & Graph & Snapshot & Tự Giám sát \\
        FLASH & Node & Snapshot & Tự Giám sát \\
        ThreaTrace & Node & Snapshot & Giám sát \\
        \hline
        HyperMamba & Event & Liên tục & Giám sát \\
        \hline
    \end{tabular}
\end{table}

\textbf{NIDS Dựa trên Nguồn gốc.} KAIROS~\cite{kairos}, MAGIC~\cite{magic}, FLASH~\cite{flash}, ThreaTrace~\cite{threatrace} không cho kết quả mức sự kiện, loại bỏ tính biến thiên đa đối tượng.

\textbf{Cơ chế Trạng thái Thời gian.} Hệ thống cũ phân rã và dùng ảnh chụp nhanh, loại bỏ trạng thái liên tục. Việc dùng SSM tích lũy theo chiều dài thời gian (Mục~\ref{sec:ablation_state}) là sáng kiến quyết định.

\textbf{Sự cân bằng trong Giám sát.} Các cách làm tự giám sát (self-supervised) đánh đổi tính chính xác sự kiện để không cần gán nhãn. Cấu trúc có giám sát của chúng tôi giúp chỉ định nhãn với độ chính xác sự kiện cao, dĩ nhiên yêu cầu dữ liệu nhãn trong huấn luyện.

\textbf{Hệ thống Phát hiện Siêu đồ thị.} Các cấu trúc trước đó (LOHA~\cite{loha}, HyperGAT~\cite{hypergat}, HRNN~\cite{hrnn}) chưa tiếp cận khái niệm siêu cạnh ở mức sự kiện và tính liên tục.

\textbf{Mô hình Không gian Trạng thái (SSM).} Kế thừa S4~\cite{s4}, S6~\cite{s6} và Mamba~\cite{mamba}, HyperMamba áp dụng mô hình cho một mạng lưới đồ thị và thực thể thay vì chỉ là token như ngôn ngữ.

\section{Kết luận}
\label{sec:conclusion}

Chúng tôi đã trình bày HyperMamba, một NIDS dựa trên nguồn gốc giúp phát hiện sự kiện tấn công xuyên tiến trình riêng biệt qua luồng giám sát thời gian thực bằng Mô hình Không gian Trạng thái Chọn lọc (SSM). Việc cắt bỏ (ablation) chứng minh đây là cơ chế then chốt: test AUPRC sụp đổ từ $0.73$ xuống $0.33$ khi thiếu nó. Trên DARPA TC E3 Theia, mô hình đạt cấp sự kiện AUPRC $0.86$ và $0.72$ ở đánh giá chéo chiến dịch --- thành công nhờ điều biến danh tính tiến trình. Với FPR cực nhỏ và tốc độ xử lý $194\times$ mức thu thập thực tế, mô hình giới hạn ở việc cần chứng minh thêm trên hệ sinh thái phần mềm lớn hơn.

\end{document}
